% corpus_supplement: Erdős problem #993
% source: https://www.erdosproblems.com/latex/993
% fetched_at: 2026-07-14T18:48:38Z
% payload_sha256: 8abb8586bb2b3cbb7949ee7160cb3aa0aa0870e44782cfae8d9972d51371360c
% statement/remarks/references extracted by erdos_ingest.extract_source_page;
% rendered by erdos_ingest.render_tex — do not edit
\documentclass{article}
\usepackage{amsmath, amssymb, amsthm}
\usepackage{hyperref}
\usepackage{geometry}
\geometry{margin=1in}

\title{Erd\H{o}s Problem \#993}
\date{Last updated: 2025-09-07}
\author{Source: erdosproblems.com}

\begin{document}
\maketitle

\noindent\textbf{Status:} OPEN \quad \textbf{No prize}

\noindent\textbf{Tags:} graph theory

\medskip
\noindent\textbf{Problem Statement:}

The independent set sequence of any tree or forest is unimodal.

In other words, if $i_k(G)$ counts the number of independent sets of vertices of size $k$ in a graph $G$, and $T$ is any tree or forest, then for some $m\geq 0$
$$i_{0}(T)\leq i_{1}(T)\leq\cdots\leq i_{m}(T)\geq i_{m+1}(T)\geq i_{m+2}(T)\geq\cdots.$$

\bigskip
\noindent\textbf{Remarks:}

A question of Alavi, Malde, Schwenk, and Erd\H{o}s \cite{AMSE87}, who showed that this is false for general graphs $G$ (in fact every possible pattern of inequalities is achieved by some graph).

The sequence which counts the number of independent sets of edges of a given size was proved to be unimodal (for any graph) by Schwenk \cite{Sc81}. In \cite{AMSE87} they also ask whether every possible unimodal pattern of inequalities is achieved by some graph.

\bigskip
\noindent\textbf{References:}

[AMSE87] Alavi, Yousef and Malde, Paresh J. and Schwenk, Allen J. and
Erd\H{o}s, Paul, The vertex independence sequence of a graph is not
constrained. Congr. Numer. (1987), 15--23.
 
 [Sc81] Schwenk, Allen J., On unimodal sequences of graphical invariants. J. Combin. Theory Ser. B (1981), 247--250.

\bigskip
\noindent\small{Source: \url{https://www.erdosproblems.com/993}}
\end{document}
